% !PRINTED 159
\weberpage{159}

\begin{center}
\S.\,49.

\gesperrt{Elimination. Theorem von Bezout.}
\end{center}

Wenn in den beiden Functionen

\weq{1}{
\begin{aligned}
f(x) &= a_0 x^n + a_1 x^{n-1} + \cdots + a_n \\
\varphi(x) &= b_0 x^m + b_1 x^{m-1} + \cdots + b_m
\end{aligned}
}

die Coëfficienten $a$, $b$ selbst wieder ganze rationale Functionen
einer Grösse $y$ sind, so wird auch die Resultante eine ganze
rationale Function von $y$ werden, die wir mit $F(y)$ bezeichnen
wollen. Die Bildung dieser Gleichung ist die Elimination von
$x$ aus den beiden Gleichungen (1).

Die Gleichung

\weq{2}{F(y) = 0}

hat alle die Werthe von $y$ zu Wurzeln, wofür die beiden
Gleichungen

\weq{3}{f(x) = 0, \qquad \varphi(x) = 0}

eine gemeinsame Wurzel haben. Ist die Bedingung (2) befriedigt,
so findet man die Werthe von $x$, die den zwei Gleichungen
(3) genügen, indem man den grössten gemeinschaftlichen Theiler
von $f(x)$ und $\varphi(x)$ aufsucht. Dieser gemeinschaftliche Theiler
ist vom ersten Grade und bestimmt also $x$ rational, wenn
nur eine solche gemeinschaftliche Wurzel vorhanden ist; im anderen
 Falle wird er von entsprechend höherem Grade, und die
Bestimmung von $x$ erfordert noch die Lösung einer Gleichung
höheren Grades.

Gehört zu jeder Wurzel der Resultante nur eine gemeinschaftliche
 Wurzel $x$, so ist die Zahl der Werthepaare $x$, $y$, die
den Gleichungen (3) genügen, gleich dem Grade der Resultante
in Bezug auf $y$. Sind also die Coëfficienten $a$ vom Grade $\nu$, $b$
vom Grade $\mu$ in Bezug auf $y$, so ist nach der Gradbestimmung
des vorigen Paragraphen der Grad der Resultante

% [source-original mathematical slip, confirmed 2026-08-25 against scan p. 184 / printed p. 159]
% The printed expression n*nu + m*mu reverses the bidegrees: the resultant has degree m in
% the a-coefficients (each of y-degree nu) and degree n in the b-coefficients (each of y-degree
% mu), so the consistent degree is m*nu + n*mu. Preserved verbatim.
\weq{5}{n\nu + m\mu;} % [sic] equation numbering jumps from (3) to (5) in print; no (4) occurs here

so gross ist also die Anzahl der gemeinsamen Wurzelpaare.
Diese Zahl kann aber noch dadurch modificirt werden, dass
möglicherweise die Coëfficienten so beschaffen sein können, dass
%CONT>>
